% ============================================================
%  Apéndice A — Instalación y puesta en marcha
%  TT 2026-B066 | ESCOM-IPN
%  Fuente: inventario del sistema (docker-compose.yml, portable/)
% ============================================================

\chapter{Instalación y puesta en marcha del sistema}
\label{ap:instalacion}

Este apéndice describe cómo se instala el sistema en la computadora del laboratorio
y qué contiene cada parte. La instalación usa Docker Compose, que levanta los cuatro
servicios del sistema con un solo archivo de configuración. No se descarga nada de
internet: las imágenes de los servicios vienen en un archivo del paquete.

\section{Requisitos}

\begin{table}[H]
\centering
\caption{Requisitos de la computadora del laboratorio.}
\label{tab:ap_requisitos}
\begin{tabular}{|p{3.6cm}|p{10.2cm}|}
\hline
\textbf{Concepto} & \textbf{Requisito} \\
\hline
Sistema operativo & Windows 10 u 11. \\
\hline
Software & Docker Desktop, que incluye Docker Compose. Conviene activar que se
inicie junto con la sesión. \\
\hline
Memoria & 16~GB en la computadora y al menos 8~GB asignados a Docker Desktop. El
servicio de análisis admite hasta 5~GB y un análisis usa cerca de 3.5~GB. \\
\hline
Disco & 20~GB libres para empezar. Las imágenes ocupan cerca de 5~GB; cada video
subido, cerca de 1~GB, y unos 200~MB más al analizarse. \\
\hline
Tarjeta gráfica & No se usa. El servicio de análisis trabaja con el procesador. \\
\hline
Internet & No hace falta para instalar ni para usar el sistema; solo para enviar
correos. \\
\hline
Energía & La computadora no debe suspenderse mientras el sistema esté en uso. \\
\hline
\end{tabular}
\end{table}

Un video de unos 6 minutos tarda entre 40 minutos y más de una hora en analizarse,
porque el análisis corre en el procesador.

\section{Servicios}

El archivo \texttt{docker-compose.yml} define cuatro servicios. Solo el último
publica un puerto: los demás se comunican dentro de la red interna de Docker y no se
alcanzan desde fuera.

\begin{table}[H]
\centering
\caption{Servicios del sistema.}
\label{tab:ap_servicios}
\begin{tabular}{|p{1.8cm}|p{9.4cm}|p{2.6cm}|}
\hline
\textbf{Servicio} & \textbf{Función} & \textbf{Puerto publicado} \\
\hline
\texttt{db} & Base de datos PostgreSQL. El esquema se aplica solo la primera vez,
con el volumen vacío. & Ninguno \\
\hline
\texttt{api} & Servidor del sistema (Spring Boot). Atiende las peticiones de la
página y corre sin privilegios de administrador. & Ninguno \\
\hline
\texttt{worker} & Toma de la base un análisis a la vez, ejecuta la detección y el
clasificador y guarda los resultados. También hace las tareas diarias: respaldo de
la base y avisos de espacio en disco. & Ninguno \\
\hline
\texttt{web} & La página del sistema, servida por nginx, que reenvía las
peticiones de la ruta \texttt{/api/} al servidor. & 8088 por omisión \\
\hline
\end{tabular}
\end{table}

Los servicios arrancan en orden: el servidor y el servicio de análisis esperan a que
la base responda, y la página espera al servidor. Los registros de cada servicio se
rotan para que no llenen el disco.

\begin{table}[H]
\centering
\caption{Volúmenes y carpetas montadas.}
\label{tab:ap_volumenes}
\begin{tabular}{|p{2.2cm}|p{4.4cm}|p{7.2cm}|}
\hline
\textbf{Servicio} & \textbf{Montaje} & \textbf{Contenido} \\
\hline
\texttt{db} & Volumen \texttt{pgdata} & Los datos de la base. \\
\hline
\texttt{db} & Archivo del esquema, de solo lectura & El esquema de la base. \\
\hline
\texttt{api} y \texttt{worker} & Volumen \texttt{videos\_reportes} & Videos, copias
para revisión, archivos de cada análisis y reportes guardados. Es el mismo volumen
para los dos servicios, así se pasan los archivos. \\
\hline
\texttt{worker} & Carpeta \texttt{modelos}, de solo lectura & El modelo del
clasificador. El servicio solo lo lee; no va dentro de la imagen. \\
\hline
\texttt{worker} & Carpeta \texttt{respaldos} & Respaldos diarios de la base, en una
carpeta de la computadora fuera de Docker. \\
\hline
\end{tabular}
\end{table}

\section{Pasos de instalación}

El paquete se entrega en dos archivos de la misma versión: un \texttt{.zip} pequeño
con la configuración, los programas de arranque, el modelo y el esquema, y un
\texttt{.tar} con las imágenes de Docker.

\begin{enumerate}
\item Descomprimir el \texttt{.zip} en una carpeta, por ejemplo \texttt{C:\textbackslash fst-portable}.
\item Copiar el \texttt{.tar} de imágenes en esa carpeta, junto a \texttt{docker-compose.yml}.
\item Abrir Docker Desktop y esperar a que indique que está corriendo.
\item Dar doble clic a \texttt{iniciar.bat}, dentro de la carpeta \texttt{portable}.
La primera vez crea el archivo \texttt{.env} a partir de una plantilla, lo abre y se
detiene.
\item Llenar en el \texttt{.env} los datos del primer administrador (variables
\texttt{ADMIN\_*}) y, si se quiere enviar correos, los de \texttt{SMTP\_*}. La llave de
sesiones y la clave de la base dicen \texttt{GENERAR} y no se tocan: \texttt{iniciar.bat}
las genera solo. Guardar y cerrar.
\item Dar doble clic otra vez a \texttt{iniciar.bat}. El programa comprueba Docker y los
datos del administrador, genera la llave y la clave, comprueba que exista el modelo,
carga las imágenes y levanta los servicios.
\item Cuando la ventana indica que el sistema está listo, el navegador abre
\texttt{http://localhost:8088}, que es la pantalla de inicio de sesión.
\end{enumerate}

Para entrar desde otras computadoras del laboratorio se ejecuta una vez
\texttt{abrir\_red.bat} como administrador, que abre el puerto solo para redes de tipo
privado, y se usa la dirección \texttt{http://NOMBRE-DE-LA-PC:8088}. La página va sin
cifrado, porque está pensada para la red del laboratorio. Para apagar el sistema se usa
\texttt{detener.bat}, que no borra datos, y para comprobar el correo,
\texttt{probar\_correo.bat}.

El compose no se ejecuta a mano: necesita datos que prepara \texttt{iniciar.bat}.

\section{Variables de entorno}

Se escriben en el archivo \texttt{.env}, sin comillas y sin espacios al final.

\begin{longtable}{|p{4.0cm}|p{9.8cm}|}
\caption{Variables del archivo \texttt{.env}.}
\label{tab:ap_variables}\\
\hline
\textbf{Variable} & \textbf{Para qué sirve} \\
\hline
\endfirsthead
\caption[]{Variables del archivo \texttt{.env} (continuación).}\\
\hline
\textbf{Variable} & \textbf{Para qué sirve} \\
\hline
\endhead
\texttt{POSTGRES\_USER}, \texttt{POSTGRES\_DB} & Usuario y nombre de la base. \\
\hline
\texttt{POSTGRES\_PASSWORD} & Clave de la base. Se genera sola en una instalación
nueva y no se cambia después de crear los datos. \\
\hline
\texttt{JWT\_SECRET} & Llave con la que se firman las sesiones. Se genera sola la
primera vez. \\
\hline
\texttt{JWT\_EXPIRES\_HOURS} & Horas que dura una sesión abierta. \\
\hline
\texttt{ADMIN\_ID} & Boleta o número de empleado del primer administrador; máximo 10
caracteres. \\
\hline
\texttt{ADMIN\_NOMBRE}, \texttt{ADMIN\_APELLIDOS} & Nombre y apellidos del primer
administrador. \\
\hline
\texttt{ADMIN\_CORREO} & Correo del primer administrador; es con el que entra. \\
\hline
\texttt{ADMIN\_PASSWORD} & Contraseña temporal, de al menos 8 caracteres. El sistema
pide cambiarla al entrar. \\
\hline
\texttt{MAX\_UPLOAD\_MB} & Tamaño máximo de un video en megabytes (2048). \\
\hline
\texttt{PUERTO\_WEB} & Puerto con el que se abre el sistema (8088 por omisión). \\
\hline
\texttt{DIRECCION\_WEB} & Dirección completa que va en los correos de recuperar
contraseña y de primer acceso. Si está vacía, se usa el nombre de la computadora. \\
\hline
\texttt{WORKER\_TZ} & Zona horaria de las tareas diarias (Ciudad de México por
omisión). \\
\hline
\texttt{MODO\_DESARROLLADOR} & Con \texttt{true} muestra herramientas de desarrollo,
como el cronómetro de las etapas. Con \texttt{false}, uso normal. \\
\hline
\texttt{EQUIPO\_NOMBRE} & Nombre opcional de la computadora para los reportes de
tiempos. \\
\hline
\texttt{SMTP\_HOST}, \texttt{SMTP\_PORT}, \texttt{SMTP\_USER}, \texttt{SMTP\_PASSWORD},
\texttt{SMTP\_AUTH}, \texttt{SMTP\_STARTTLS}, \texttt{SMTP\_FROM} & Datos del servidor
de correo. Si \texttt{SMTP\_HOST} está vacío no se envía ningún correo y los avisos
salen solo en el panel de notificaciones. \\
\hline
\end{longtable}

El compose fija por su cuenta otras variables, como la dirección de la base, el puerto
interno del servidor y la ruta del modelo. Otras tienen valor por omisión y no hace
falta escribirlas: los minutos que vale el enlace para ver un video (120) y el de
recuperar contraseña (60), la hora de las tareas diarias (3), las horas máximas que un
análisis puede quedar en pausa (72) y los porcentajes de disco para avisar (80) y para
limpiar (90).

\section{Primer administrador}

El primer administrador no se crea con una pantalla ni con un comando. Lo crea el
servidor al arrancar, con los datos de las variables \texttt{ADMIN\_*}.

Si falta alguno de los cinco datos, o si el identificador pasa de 10 caracteres, el
servidor no crea nada y avisa cuál falta. Si el administrador ya existe, no lo vuelve
a crear. El usuario y su registro de administrador se crean en una sola operación, de
modo que nunca queda un usuario sin su registro. En el primer ingreso el sistema pide
cambiar la contraseña temporal. Las demás cuentas solo las crea un administrador, y
si alguien olvida su contraseña, un administrador puede darle una temporal.

\section{Instalación para desarrollo}

Quien arma el paquete o modifica el sistema trabaja desde el repositorio. Coloca el
modelo en la carpeta \texttt{modelos} y levanta el entorno de desarrollo con
\texttt{docker compose --profile dev up --build} dentro de la carpeta del servidor. Ese
entorno publica el puerto 5433 para la base y el 8084 para el servidor, activa la
documentación de la interfaz del servidor y dirige los correos a un buzón de prueba que
se consulta en el navegador. El paquete de instalación se arma con el programa
\texttt{empaquetar.ps1}, que construye las imágenes y genera el \texttt{.tar} y el
\texttt{.zip}.
